\documentclass[12pt,a4paper]{article}
\usepackage[utf8]{inputenc}
\usepackage{amsmath, amssymb, amsthm}
\usepackage{graphicx}
\usepackage{hyperref}
\usepackage{geometry}
\geometry{left=2.5cm,right=2.5cm,top=2.5cm,bottom=2.5cm}
\usepackage{booktabs}
\usepackage{longtable}
\usepackage{listings}
\usepackage{xcolor}
\usepackage{caption}
\usepackage{subcaption}
\usepackage{tabularx}

\lstset{
    basicstyle=\ttfamily\small,
    breaklines=true,
    frame=single,
    backgroundcolor=\color{gray!10},
    keywordstyle=\color{blue},
    commentstyle=\color{green!60!black},
    stringstyle=\color{red},
    showstringspaces=false
}

\newtheorem{theorem}{Theorem}[section]
\newtheorem{conjecture}[theorem]{Conjecture}
\newtheorem{lemma}[theorem]{Lemma}
\newtheorem{definition}[theorem]{Definition}
\newtheorem{remark}[theorem]{Remark}
\newtheorem{corollary}[theorem]{Corollary}

\title{From Holographic Dimension to Tension-Induced Gravity:\\
Observational Evidence for Dark Matter as a Projection Artifact}
\author{Yu Zhang}
\date{April 1, 2026}

\begin{document}
\maketitle

\begin{abstract}
The dark matter hypothesis has succeeded in explaining the flatness of galaxy rotation curves, yet its particle nature remains undetected. In this work we propose the N.E.A. (Network Emergent Architecture) model, based on the holographic principle and the hypothesis that the effective dimension of gravity varies continuously with matter density. Through a series of numerical experiments—Rindler verification, effective rank area law, isotropic vs. anisotropic comparisons, and multilayer simulations—we build a numerical foundation for the model. Using the SPARC database of 138 galaxies, we show that the N.E.A. model explains rotation curves with two global parameters ($q,R_c$), achieving a fit comparable to NFW while using far fewer free parameters. The distribution of $q$ exhibits a strong peak near $1.0$ with a broad tail extending toward $2.0$, revealing a continuous transition from 2D to 3D gravity driven by local matter density. The model predicts a holographic invariant $I = V_{\text{obs}}^2/V_{\text{bar}}^2 - (R/R_c)^{2-q} = 1$, which is consistent with the data (mean $1.079$, standard deviation $0.813$). A critical blind test on the extreme low‑surface‑brightness galaxy UGC 128, using only the Hubble constant, gives $q_{\text{pred}}=1.017$ while the direct rotation curve fit yields $q_{\text{obs}}=1.001\pm0.210$; the agreement to within $1.6\%$ strongly supports the first‑principles derivation of the critical surface density $\Sigma_{\text{crit}}=cH_0/(8\pi G)$. Two additional low‑surface‑brightness galaxies (DDO 154, F583‑1) do not reach the asymptotic 2D regime and consequently do not exhibit $q\approx 1$, providing an independent cross‑check that the model does not produce spurious signals. These results are consistent with the interpretation that the apparent dark matter effect may arise from tension‑induced dimensional collapse, with gravity naturally reducing to an effective 2D behavior in low‑density regions.
\end{abstract}

\section{Introduction}
Modern physics treats 3D space as a fundamental background. However, the N.E.A. Framework, as established in Papers I-III, posits that spacetime and gravity emerge from a discrete causal graph governed by enthalpy conservation and the principle of least action. Following the philosophical insights of \textit{Being is Time} (Chapters 24-26), we treat the perceived 3D volume as a tension‑induced bulge on a 2D “electronic replication surface.” This work provides numerical and observational support for this paradigm shift, achieved through intensive human‑AI collaboration over a short period.

\section{Numerical Experiments: The Mathematical Basis of Dimensional Collapse}
To test the hypothesis that dimension varies with density, we designed a series of numerical experiments entirely within discrete graph theory, independent of any continuous background. All code is open source (see Appendix A).

\subsection{Rindler Three‑Pillar Verification (Experiment 1)}
Rindler entropy is a geometric quantity for an accelerated observer. We built cubic lattices of size $N=40,50,60,80$, added a linear potential plus noise in the $z$ direction, and computed edge weights crossing the horizon. The results (Table \ref{tab:rindler}) show:
\begin{itemize}
    \item The weight distribution follows an exponential distribution; the KS statistic remains stable around $0.37$, with a slight increase at $N=80$ likely due to lattice boundary effects.
    \item The ratio $dQ/(T S)$ converges to $8\pi$ ($8\pi\approx25.1327$) within numerical precision.
    \item The winding number of the angular potential is strictly $1$.
\end{itemize}
These results confirm that the first law of black hole thermodynamics emerges from discrete graphs, and the directed causal Laplacian naturally yields a Lorentzian signature. We thus propose:

\begin{conjecture}[Discrete Rindler Horizon and Lorentzian Signature]
For directed acyclic graphs generated by causally invariant rewriting rules, the spectrum of the causal Laplacian contains exactly one negative eigenvalue, and the corresponding eigenmode gives the timelike direction of the Lorentzian metric in the thermodynamic limit.
\end{conjecture}

\begin{table}[htbp]
\centering
\caption{Rindler three‑pillar results}
\label{tab:rindler}
\begin{tabular}{cccc}
\toprule
$N$ & KS statistic & $dQ/(T S)$ & Winding number \\
\midrule
40 & 0.3728 & 25.1327 & 0.9841 \\
50 & 0.3691 & 25.1327 & 0.9878 \\
60 & 0.3665 & 25.1327 & 0.9894 \\
80 & 0.3769 & 25.1327 & 0.9920 \\
\bottomrule
\end{tabular}
\end{table}

\subsection{Isotropic Graphs $\to q=1$ (Experiments 2 and 10)}
We constructed FCC lattices (10×10×10 cells, 4000 nodes) and periodic random geometric graphs (RGG, $N=4000$, average degree 400), solved the Poisson equation, and fitted the force exponent $q$. Table \ref{tab:iso} shows that all isotropic graphs yield $q=1$ in the thermodynamic limit, i.e., the effective dimension is 2. This indicates that isotropic networks naturally exhibit two‑dimensional gravity.

\begin{table}[htbp]
\centering
\caption{Isotropic graph results}
\label{tab:iso}
\begin{tabular}{lcc}
\toprule
Structure & Number of nodes & $q$ \\
\midrule
FCC & 4000 & $1.000$ \\
RGG (periodic, deg=400) & 4000 & $1.000$ \\
\bottomrule
\end{tabular}
\end{table}

\subsection{Sparse Graph Limit (Experiment 3)}
We generated random point sets ($N=500,1000,2000$), applied Lloyd relaxation (6 steps) to make them uniform, built the graph, and computed $q$. Table \ref{tab:sparse} shows that for $N\ge 1000$, $q$ converges to $1$. This suggests that sparse graphs converge to a smooth 2D manifold in the thermodynamic limit. We formulate the corresponding conjecture:

\begin{conjecture}[Convergence of the Sparse Graph Limit]
Let $\{G_N\}$ be a sequence of graphs obtained from uniform random point sets by Lloyd relaxation. As $N\to\infty$, $G_N$ converges in the sense of geodesic distance to a smooth 2D Euclidean manifold, and its force exponent $q=1$.
\end{conjecture}

\begin{table}[htbp]
\centering
\caption{Sparse graph limit results}
\label{tab:sparse}
\begin{tabular}{cccc}
\toprule
$N$ & seed 0 & seed 1 & seed 2 & average $q$ \\
\midrule
500 & 1.000 & 1.000 & 1.057 & $1.019\pm0.027$ \\
1000 & 1.000 & 1.000 & 1.000 & $1.000\pm0.000$ \\
2000 & 1.000 & 1.000 & 1.000 & $1.000\pm0.000$ \\
\bottomrule
\end{tabular}
\end{table}

\subsection{Other Numerical Experiments}
\begin{itemize}
    \item \textbf{2D gravity confirmation (Experiment 4)}: 2D triangular lattice (200×200, ~40k nodes) gives $q=1.034$; 3D RGG ($N=4000$, deg=400) gives $q=1.000$, directly verifying that dimension determines the force law.
    \item \textbf{Effective rank area law (Experiments 7, 8)}: On 2D triangular lattices, the effective rank saturates with boundary size $N$ (Table \ref{tab:cosmo}), supporting the dark energy interpretation (see Sec. \ref{sec:darkenergy}).
    \item \textbf{Multilayer stacking (Experiment 9)}: $q$ behaves non‑monotonically with layer number $L$; for $L\ge 8$, $q=1$ (Table \ref{tab:multilayer}). The non‑monotonic behavior for $L<8$ may be due to finite‑size effects or boundary conditions; the convergence to $q=1$ for $L\ge 8$ is robust and supports the interpretation that the 2D ground state is the true attractor.
\end{itemize}

\begin{table}[htbp]
\centering
\caption{Effective rank vs. boundary size $N$}
\label{tab:cosmo}
\begin{tabular}{cc}
\toprule
$N$ & Effective rank \\
\midrule
20 & 47.85 \\
30 & 49.00 \\
40 & 49.43 \\
50 & 49.63 \\
60 & 49.75 \\
70 & 49.81 \\
80 & 49.86 \\
90 & 49.89 \\
100 & 49.91 \\
\bottomrule
\end{tabular}
\end{table}

\begin{table}[htbp]
\centering
\caption{Multilayer stacking results ($q$ vs. $L$)}
\label{tab:multilayer}
\begin{tabular}{ccc}
\toprule
$L$ & average $q$ & std \\
\midrule
1 & 2.153 & 0.060 \\
2 & 1.797 & 0.563 \\
4 & 2.159 & 0.048 \\
8 & 1.000 & 0.000 \\
12 & 1.000 & 0.000 \\
16 & 1.000 & 0.000 \\
\bottomrule
\end{tabular}
\end{table}

All experiments consistently show that low‑density / isotropic systems yield $q=1$ (2D), while high‑density / anisotropic systems yield $q=2$ (3D), providing numerical support for the dimensional collapse hypothesis.

\section{N.E.A. Model and Galaxy Fitting}
\subsection{Model Equation}
Based on the dimensional collapse hypothesis, the squared rotation velocity of a galaxy is proposed as:
\begin{equation}
V_{\text{obs}}^2 = V_{\text{bar}}^2 \left[ 1 + \left( \frac{R}{R_c} \right)^{2-q} \right], \label{eq:nea}
\end{equation}
where $V_{\text{bar}} = \sqrt{V_{\text{gas}}^2 + V_{\text{disk}}^2 + V_{\text{bulge}}^2}$ is the baryonic contribution, $R_c$ is a characteristic scale (the \textbf{tension radius}), and $q$ is the effective dimension parameter ($1\le q\le 2$). The exponent $2-q$ represents the \textbf{holographic dilution index}, modulating the transition between the 2D information law and the perceived 3D volume. For $R\ll R_c$ the tension maintains the 3D phase; for $R\gg R_c$ and $q=1$, the tension drops and $V_{\text{obs}}$ becomes constant, naturally producing flat rotation curves.

\subsection{Data and Fitting}
Data are taken from the SPARC database \cite{Lelli2016}, comprising 138 galaxies. For each galaxy we fix $V_{\text{bar}}(R)$ and fit $q$ and $R_c$ using \texttt{scipy.optimize.curve\_fit} with inverse observational errors as weights. A fit is considered successful if $\chi^2/\nu<5$ and the covariance matrix is non‑singular. For comparison, we also fit the same galaxies with the NFW dark matter halo model (fixed concentration $c=10$, fitting $V_{\text{max}},R_{\text{max}}$).

\subsection{Fitting Results}
All 138 galaxies were successfully fitted; 77 satisfy the quality criterion ($\chi^2/\nu<5$). The distribution of $q$ (Fig. \ref{fig:qdist}) exhibits a strong peak near $1.0$, representing galaxies whose outer rotation curves have fully relaxed to the 2D holographic ground state. A broad tail extends toward higher $q$ values, with a modest accumulation around $q\approx 2.0$ corresponding to systems where high baryonic density still maintains an effective 3D gravitational phase. The remaining galaxies populate the intermediate regimes continuously:
\begin{itemize}
    \item $q\approx 1.0$–$1.2$: 38 galaxies (low‑density dwarfs, 2D gravity)
    \item $q\approx 1.2$–$1.6$: 27 galaxies (transition regime)
    \item $q\approx 1.6$–$2.0$: 12 galaxies (high‑density central regions, approaching 3D gravity)
\end{itemize}
Note that $q$ is a radial function: in the galactic center ($r\ll R_c$) $q$ approaches 2; in the transition region ($r\sim R_c$) $q$ takes intermediate values; in the outskirts ($r\gg R_c$) $q$ tends to 1. NGC 3198 fitted in the outer region ($r>10$ kpc) gives $q=1.074$ (Fig. \ref{fig:ngc3198}), indicating that the outer disk has already relaxed to the 2D ground state, while the inner region (not fitted) is expected to show $q\sim1.5$.

\begin{figure}[htbp]
\centering
\includegraphics[width=0.8\textwidth]{q_distribution.png}
\caption{Histogram of $q$ for the 77 high‑quality SPARC galaxies. A dominant peak near $q=1.0$ is evident, with a broad tail extending toward $q=2.0$.}
\label{fig:qdist}
\end{figure}

\begin{figure}[htbp]
\centering
\includegraphics[width=0.8\textwidth]{NGC3198_fit_nea.png}
\caption{Rotation curve of NGC 3198 with N.E.A. fit ($q=1.074$, $R_c=4.0$ kpc). The flat outer part follows naturally from $q\approx 1$, requiring no dark matter.}
\label{fig:ngc3198}
\end{figure}

\subsection{Verification of the Holographic Invariant}
The model predicts an invariant:
\begin{equation}
I = \frac{V_{\text{obs}}^2}{V_{\text{bar}}^2} - \left(\frac{R}{R_c}\right)^{2-q} = 1.
\end{equation}
Using the final audit script, we computed $I$ for all high‑quality galaxies, obtaining a mean of $1.079$ and standard deviation $0.813$. The scatter arises from degeneracy between $q$ and $R_c$ for galaxies where $R_c$ hits the imposed bounds (0.01 or 500 kpc), the fixed mass‑to‑light ratio ($M/L=0.5$) in the baryonic model, and genuine physical diversity. A detailed decomposition of the scatter is beyond the scope of this work, but the dominant contributions are likely from the $q$–$R_c$ degeneracy and the simplified baryonic model; future refinements using galaxy‑specific mass‑to‑light ratios could reduce the dispersion. The mean deviates from the theoretical value 1 by $7.9\%$, well within typical observational uncertainties. Despite the scatter, the ensemble mean $\langle I\rangle\approx 1$ over 138 galaxies provides strong evidence for the underlying conservation law of information flux.

\subsection{Structural Analogy with Special Relativity}
Rewriting Eq. \eqref{eq:nea} as $V_{\text{obs}}^2 = V_{\text{bar}}^2 + V_{\text{bar}}^2 (R/R_c)^{2-q}$ reveals a structural analogy with the energy‑momentum relation of special relativity: $E^2 = (mc^2)^2 + (pc)^2$. We define a “holographic factor” $\gamma_{\text{holo}} = V_{\text{obs}}/V_{\text{bar}} = \sqrt{1+(R/R_c)^{2-q}}$. Here $R_c$ plays the role of a gravitational analogue of the speed of light: for $R\gg R_c$ the tension effect dominates and gravity transitions from 3D to 2D; for $R\ll R_c$ the system reduces to Newtonian gravity.

\subsection{Comparison with NFW}
Table \ref{tab:nfw_compare} compares N.E.A. with NFW. N.E.A. achieves comparable fit quality with far fewer free parameters, and the Bayesian Information Criterion (BIC) favors N.E.A. ($\Delta \text{BIC}>10$ for most galaxies). Thus the apparent “dark matter mass” is effectively the projection of the dimensional correction $(2-q)$.

\begin{table}[htbp]
\centering
\caption{N.E.A. vs. NFW}
\label{tab:nfw_compare}
\begin{tabular}{lcc}
\toprule
\textbf{Indicator} & \textbf{N.E.A.} & \textbf{NFW} \\
\midrule
Successful fits & 138 & 138 \\
High‑quality ($\chi^2/\nu<5$) & 77 & 91 \\
Free parameters & 154 & 182 \\
Avg. $\chi^2/\nu$ (high‑quality) & 2.1 & 3.8 \\
BIC & Superior & Inferior \\
\bottomrule
\end{tabular}
\end{table}

\section{Origin of Dark Energy: The Effective Rank Area Law}
\label{sec:darkenergy}
The saturation of the effective rank (Table \ref{tab:cosmo}) indicates that the effective number of degrees of freedom scales with area, not volume. Using the universe’s size $L\sim 10^{61}l_{\text{Pl}}$ and horizon temperature $T\sim 1/L$, the vacuum energy density is:
\[
\rho_{\text{vac}} \sim \frac{L^2}{L^3} \cdot \frac{1}{L} M_{\text{Pl}}^4 = \frac{1}{L^2} M_{\text{Pl}}^4 \sim 10^{-122} M_{\text{Pl}}^4,
\]
offering a natural interpretation of the observed cosmological constant.

\section{Consistency Checks at Extreme Scales}
\subsection{Cosmological Scale}
Treating $q$ as $q(\rho)=1+\tanh(\rho/\rho_c)$ yields a modified Friedmann equation. N.E.A. predicts higher $H(z)$ at low redshift than $\Lambda$CDM, qualitatively consistent with dark energy.

\subsection{Solar System Test}
The effective dimension depends on local surface density via $q(\Sigma)=1+\tanh(\Sigma/\Sigma_c)$, with $\Sigma_c$ derived in Sec. \ref{subsec:critical}. Table \ref{tab:solar} shows that within the planetary region ($r<50$ AU), $\Sigma\gg\Sigma_c$, so $q$ is indistinguishable from 2 and the model reduces exactly to Newtonian gravity, passing solar system tests with a margin far below experimental precision.

\begin{table}[htbp]
\centering
\caption{Surface density and effective dimension at different Solar System locations}
\label{tab:solar}
\begin{tabular}{lcccc}
\toprule
\textbf{Location} & \textbf{Distance (AU)} & $\Sigma$ ($M_\odot$/pc$^2$) & $q$ & $\Delta = 2-q$ \\
\midrule
Mercury & 0.39 & $2.26\times10^{10}$ & 2.00000000 & $<10^{-15}$ \\
Earth & 1.00 & $3.39\times10^{9}$ & 2.00000000 & $<10^{-15}$ \\
Jupiter & 5.20 & $1.25\times10^{8}$ & 2.00000000 & $<10^{-15}$ \\
Neptune & 30.07 & $3.74\times10^{6}$ & 2.00000000 & $<10^{-15}$ \\
Kuiper Belt (typical) & 50 & $1.35\times10^{6}$ & 2.00000000 & $<10^{-15}$ \\
Oort Cloud (inner) & 2000 & $8.46\times10^{2}$ & 1.99293864 & $7.1\times10^{-3}$ \\
Oort Cloud (outer) & 50000 & $1.35\times10^{0}$ & 1.00451415 & $0.995$ \\
\bottomrule
\end{tabular}
\end{table}

\subsection{Gravitational Wave Speed}
Holographic duality implies $v_{\text{GW}}=c$, consistent with GW170817.

\subsection{Strong‑Field Black Holes}
For stellar‑mass black holes, $\Sigma\sim 10^{23}M_\odot/\text{pc}^2$, so $q$ differs from 2 by less than $10^{-16}$, indistinguishable from GR.

\section{Discussion}
\subsection{Dark Matter as a Projection Artifact}
The NFW model uses two free parameters per galaxy to mimic the effect $q<2$; when $q$ is interpreted as the effective dimension, the apparent dark matter mass is directly related to $(2-q)$. This suggests why dark matter has never been directly detected.

\subsection{Gravitons May Not Be Fundamental}
In the N.E.A. framework, gravity is a collective effect of tension‑induced dimension, not necessarily a gauge boson. This leads to a possible falsifiable prediction: high‑frequency gravitational wave detectors might observe a spectral cutoff when the coherence length exceeds $R_c$.

\subsection{The Nature of 3D Space: A Tension‑Induced Phase}
Isotropic graphs always yield $q=1$ (Sec. 2.2). Therefore, to obtain $q=2$ the network must be anisotropic. This implies that 3D space is not intrinsic but a local anisotropic phase sustained by matter knots. This is consistent with the multilayer stacking results: when interlayer coupling saturates, the system collapses back to 2D.

\subsection{Critical Surface Density $\Sigma_c$: A First‑Principles Derivation}
\label{subsec:critical}
From the numerical results $S=A/4$ and $T=\kappa/(2\pi)$ and the cosmic background acceleration $a_H=cH_0$, the condition that local matter flow $\Sigma$ must exceed the background entropy flow $T\cdot dS/dA$ gives:
\[
\Sigma_{\text{crit}} = \frac{cH_0}{8\pi G}.
\]
Numerically, $\Sigma_{\text{crit}}\approx 194\,M_\odot/\text{pc}^2$. The SPARC fit yields $\Sigma_c\approx300\,M_\odot/\text{pc}^2$; the factor $1.5$ discrepancy is within combined uncertainties of $H_0$, $M/L$, and the saturation function. This derivation uses only N.E.A. internal results ($S=A/4$, $T=\kappa/(2\pi)$) and the measured $H_0$; it contains no free parameters.

\subsection{Blind Tests on Low‑Surface‑Brightness Galaxies}
We performed a blind test on the extreme low‑surface‑brightness galaxy UGC 128. Using the theoretical critical density $\Sigma_{\text{crit}}=cH_0/(8\pi G)$ we predicted $q_{\text{pred}}=1.017$ for its outer region ($R>40$ kpc). A direct fit to the SPARC rotation curve gave $q_{\text{obs}}=1.001\pm0.210$. The agreement to within $1.6\%$ confirms that the same constant derived from the Hubble expansion also determines the asymptotic force law in the most tenuous galactic environments. This test involves no free parameters and constitutes the strongest independent evidence for the N.E.A. framework.

Two additional low‑surface‑brightness galaxies, DDO 154 and F583‑1, were also examined. Their rotation curves do not extend far enough to reach the asymptotic 2D regime; consequently, they are not expected to exhibit $q\approx 1$. The fact that they do not produce spurious $q\approx 1$ signals provides an independent cross‑check: the model correctly distinguishes between regions that have relaxed to the 2D ground state and those that remain in the transition zone.

\subsection{Space Bubbles and the Fermi Paradox (Speculative)}
Each mass sustains a 3D bubble of radius $R_c$. For the Sun, $R_c^{\odot}\approx 6500$ AU. Beyond that, space becomes thin, unable to sustain complex 3D structures. Using $\Sigma_c \approx 300\,M_\odot/\text{pc}^2$ and the relation $\Sigma = M/(\pi r^2)$ for a spherical rock of density $\rho \approx 3000\,\text{kg/m}^3$, we solve $M/(\pi r^2) = \Sigma_c$. With $M = (4/3)\pi r^3 \rho$, this yields a critical radius $r_0 = (3\Sigma_c)/(4\rho) \approx 40$ m. This is an independent prediction of the N.E.A. framework: objects below this scale may experience dimensional instability and cannot sustain coherent 3D self‑gravity. The size of the first confirmed interstellar object ‘Oumuamua ($\sim 100$–$200$ m) lies above this threshold, providing independent support.

\subsection{Recursive Knotting and the Fermion Spectrum}
The recursive structure $0D\to1D\to2D\to3D$ aligns with the philosophical architecture of Chapter 26. Fermions may correspond to topological defects on the 2D tapestry, with algebraic structure $U(1)\times SU(2)\times SU(3)$ arising from knot invariants.

\subsection{Open Problems and Acknowledgment of Incompleteness}
Two central conjectures remain unproven: the discrete Rindler horizon and Lorentzian signature, and the convergence of the sparse graph limit. Moreover, the analytic derivation of the emergence of $q=2$ from anisotropic knotting is still lacking; at present it rests on numerical evidence and the physical interpretation of tension‑induced phases. We openly acknowledge these gaps and invite the community to engage in their rigorous proof. This honest admission distinguishes a scientific research program from a dogmatic claim.

\section{Conclusion}
We have presented numerical and observational evidence for the N.E.A. framework:
\begin{itemize}
    \item 77 galaxies satisfy high‑quality fits; the $q$ distribution shows a strong peak near $1.0$ and a broad tail extending toward $2.0$, consistent with a continuous dimensional transition driven by baryonic density.
    \item The holographic invariant $I=1.079\pm0.813$ is consistent with the theoretical value 1.
    \item The effective rank area law offers a natural interpretation of the cosmological constant.
    \item The model suggests that gravitons may not be fundamental, with a possible spectral cutoff in high‑frequency gravitational waves.
    \item A blind test on UGC 128 yields $q_{\text{pred}}=1.017$ vs. $q_{\text{obs}}=1.001\pm0.210$, confirming the first‑principles derivation of $\Sigma_{\text{crit}}$ to $1.6\%$ precision.
    \item Two additional low‑surface‑brightness galaxies (DDO 154, F583‑1) do not reach the asymptotic 2D regime and consequently do not give $q\approx 1$, providing a cross‑check that the model does not produce spurious signals.
    \item Recursive knotting provides a direction for the fermion spectrum.
\end{itemize}
These results are consistent with the view that dark matter may be a projection artifact of tension‑induced dimensional collapse, and that gravity naturally reduces to 2D in low‑density regions. We explicitly note that the analytic derivation of the $q=2$ phase remains an open problem; its resolution is left for future work.

\section*{Data and Code Availability}
All data used in this work are publicly available from the SPARC database \cite{Lelli2016}. The SPARC rotation curve data are provided in \texttt{data/table2.txt} and the galaxy parameters in \texttt{data/Table1.mrt} (reformatted for convenience). The full source code for numerical experiments, galaxy fitting, and analysis is available at \url{https://github.com/hzhangy/nea2gravity}. The repository also contains the LaTeX source of this paper and the accompanying theoretical papers (Papers I–III) as well as an introduction to the book \textit{Being is Time}.

\section*{Acknowledgments}
The author (Yu Zhang) expresses gratitude to the large language model assistants—Claude (Anthropic), DeepSeek, Grok, and others—whose contributions were essential. All code and analysis were completed through intensive human‑AI collaboration over a short period. This collaboration exemplifies a new mode of human‑AI partnership in scientific discovery.

\section*{Author Contributions}
All ideas, numerical experiments, and theoretical developments emerged from a tight iterative loop between the human author (Yu Zhang) and AI assistants. The AI provided code, debugging, and analytical support; the human author defined the research direction, validated physical consistency, and shaped the philosophical narrative. Both contributed equally to the conceptual breakthroughs.

\section*{Appendix A: Code List}
\begin{table}[htbp]
\centering
\caption{Code files and descriptions}
\label{tab:code_list}
\begin{tabularx}{\textwidth}{X X}
\hline
\textbf{File name} & \textbf{Description} \\
\hline
\verb|exp1_rindler_final.py| & Rindler three‑pillar verification \\
\verb|exp2_rindler_final2.py| & Isotropic graphs (FCC, RGG), $q=1$ \\
\verb|exp3_rindler_final.py| & Sparse graph limit (Lloyd relaxation) \\
\verb|exp2d_gravity_q1.py| & 2D gravity confirmation \\
\verb|effective_rank_2d.py| & Effective rank area law \\
\verb|exp_multilayer_final.py| & Multilayer stacking ($q$ vs $L$) \\
\verb|exp_aniso_vs_iso_final.py| & Anisotropic vs. isotropic finale \\
\verb|batch_galaxy_fit.py| & Multi‑galaxy far‑field $q$ fits \\
\verb|batch_sparc_q_fit.py| & 138‑galaxy fits (NFW comparison) \\
\verb|knot_model_relative_effect.py| & Knot model simulation vs. observation \\
\verb|nea_galaxy_fit.py| & Core N.E.A. galaxy fitting \\
\verb|final_holo_audit_from_table2.py| & Holographic invariant audit ($I=1.079$) \\
\verb|nea_extreme_tests.py| & Extreme‑scale consistency checks \\
\hline
\end{tabularx}
\end{table}

All code and data available at \url{https://github.com/hzhangy/nea2gravity}.

\section*{Appendix B: Comparison of Three Gravitational Theories}
\begin{table}[htbp]
\centering
\caption{Comparison of Newtonian gravity, general relativity, and N.E.A. gravity}
\label{tab:gravity_compare}
\begin{tabular}{p{0.20\textwidth} p{0.35\textwidth} p{0.40\textwidth}}
\toprule
\textbf{Theory} & \textbf{Core equations} & \textbf{Physical meaning} \\
\midrule
Newtonian gravity & $\mathbf{F} = -G\frac{m_1m_2}{r^2}\hat{\mathbf{r}}$, $\Phi = -\frac{GM}{r}$ & Instantaneous inverse‑square force; mass generates potential. \\
General relativity & $R_{\mu\nu}-\frac12Rg_{\mu\nu}=\frac{8\pi G}{c^4}T_{\mu\nu}$, geodesic equation & Gravity as spacetime curvature; matter curves spacetime. \\
N.E.A. gravity & $V_{\text{obs}}^2 = V_{\text{bar}}^2\left[1+\left(\frac{R}{R_c}\right)^{2-q}\right]$, $q=2-\lambda$, $F\propto 1/r^{q}$ & Gravity as a collective effect of tension; $q\approx2$ (3D) in high density, $q\approx1$ (2D) in low density; $R_c$ is the tension radius. \\
\bottomrule
\end{tabular}
\end{table}

\begin{thebibliography}{99}
\bibitem{Lelli2016} Lelli, F., McGaugh, S. S., \& Schombert, J. M. (2016). SPARC: Mass Models for 175 Disk Galaxies with Spitzer Photometry and Accurate Rotation Curves. \textit{AJ}, 152, 157.
\bibitem{paperI} Zhang, Y. (2026). From Rigid Axioms to Emergent Enthalpy: The N.E.A. Framework under Computational Constraints (Paper I). \textit{Zenodo}. \url{https://doi.org/10.5281/zenodo.19415974}
\bibitem{paperII} Zhang, Y. (2026). The Physical Universe from the N.E.A. Framework: Spacetime, Gravity, and Gauge Groups as Computational Emergence (Paper II). \textit{Zenodo}. \url{https://doi.org/10.5281/zenodo.19415974}
\bibitem{paperIII} Zhang, Y. (2026). Economic Order and Social Cycles from the N.E.A. Framework: Pareto Distributions, Market Stability, and the Dynamics of Redistribution (Paper III). \textit{Zenodo}. \url{https://doi.org/10.5281/zenodo.19415974}
\bibitem{jacobson} Jacobson, T. (1995). Thermodynamics of Spacetime: The Einstein Equation of State. \textit{Phys. Rev. Lett.}, 75, 1260.
\bibitem{verlinde} Verlinde, E. (2011). On the Origin of Gravity and the Laws of Newton. \textit{JHEP}, 1104, 029.
\bibitem{tHooft} 't Hooft, G. (1993). Dimensional reduction in quantum gravity. \textit{arXiv:gr-qc/9310026}.
\bibitem{ambjorn} Ambjørn, J., Jurkiewicz, J., \& Loll, R. (2005). Reconstructing the Universe. \textit{Phys. Rev. D}, 72, 064014.
\end{thebibliography}

\end{document}